% Options for packages loaded elsewhere
\PassOptionsToPackage{unicode}{hyperref}
\PassOptionsToPackage{hyphens}{url}
\documentclass[
  ignorenonframetext,
]{beamer}
\newif\ifbibliography
\usepackage{pgfpages}
\setbeamertemplate{caption}[numbered]
\setbeamertemplate{caption label separator}{: }
\setbeamercolor{caption name}{fg=normal text.fg}
\beamertemplatenavigationsymbolsempty
% remove section numbering
\setbeamertemplate{part page}{
  \centering
  \begin{beamercolorbox}[sep=16pt,center]{part title}
    \usebeamerfont{part title}\insertpart\par
  \end{beamercolorbox}
}
\setbeamertemplate{section page}{
  \centering
  \begin{beamercolorbox}[sep=12pt,center]{section title}
    \usebeamerfont{section title}\insertsection\par
  \end{beamercolorbox}
}
\setbeamertemplate{subsection page}{
  \centering
  \begin{beamercolorbox}[sep=8pt,center]{subsection title}
    \usebeamerfont{subsection title}\insertsubsection\par
  \end{beamercolorbox}
}
% Prevent slide breaks in the middle of a paragraph
\widowpenalties 1 10000
\raggedbottom
\AtBeginPart{
  \frame{\partpage}
}
\AtBeginSection{
  \ifbibliography
  \else
    \frame{\sectionpage}
  \fi
}
\AtBeginSubsection{
  \frame{\subsectionpage}
}
\usepackage{iftex}
\ifPDFTeX
  \usepackage[T1]{fontenc}
  \usepackage[utf8]{inputenc}
  \usepackage{textcomp} % provide euro and other symbols
\else % if luatex or xetex
  \usepackage{unicode-math} % this also loads fontspec
  \defaultfontfeatures{Scale=MatchLowercase}
  \defaultfontfeatures[\rmfamily]{Ligatures=TeX,Scale=1}
\fi
\usepackage{lmodern}
\ifPDFTeX\else
  % xetex/luatex font selection
\fi
% Use upquote if available, for straight quotes in verbatim environments
\IfFileExists{upquote.sty}{\usepackage{upquote}}{}
\IfFileExists{microtype.sty}{% use microtype if available
  \usepackage[]{microtype}
  \UseMicrotypeSet[protrusion]{basicmath} % disable protrusion for tt fonts
}{}
\makeatletter
\@ifundefined{KOMAClassName}{% if non-KOMA class
  \IfFileExists{parskip.sty}{%
    \usepackage{parskip}
  }{% else
    \setlength{\parindent}{0pt}
    \setlength{\parskip}{6pt plus 2pt minus 1pt}}
}{% if KOMA class
  \KOMAoptions{parskip=half}}
\makeatother
\usepackage{longtable,booktabs,array}
\newcounter{none} % for unnumbered tables
\usepackage{calc} % for calculating minipage widths
\usepackage{caption}
% Make caption package work with longtable
\makeatletter
\def\fnum@table{\tablename~\thetable}
\makeatother
\setlength{\emergencystretch}{3em} % prevent overfull lines
\providecommand{\tightlist}{%
  \setlength{\itemsep}{0pt}\setlength{\parskip}{0pt}}
\usepackage{bookmark}
\IfFileExists{xurl.sty}{\usepackage{xurl}}{} % add URL line breaks if available
\urlstyle{same}
\hypersetup{
  pdftitle={Cognitive Integration: Active Inference and CEREBRUM},
  hidelinks,
  pdfcreator={LaTeX via pandoc}}

\title{\texorpdfstring{Cognitive Integration: Active Inference and
CEREBRUM}{Cognitive Integration: Active Inference and CEREBRUM}}
\author{\texorpdfstring{}{}}
\date{}

\begin{document}
\frame{\titlepage}

\section{Cognitive Integration: Active Inference, CEREBRUM, and
Diagrammatic Reasoning}\label{sec:cognitive-integration}

\begin{frame}{The Missing Layer: A Dynamic Process Theory of Cognition}
\protect\phantomsection\label{the-missing-layer-a-dynamic-process-theory-of-cognition}
The preceding sections have developed a rich mathematical infrastructure
for analyzing case systems---categorical, type-logical, distributional,
enriched, and topos-theoretic. But these frameworks are \emph{static}:
they describe the structure of case without explaining how a cognitive
agent \emph{uses} that structure in real-time language understanding and
production. What is needed is a \emph{process theory} that explains how
case-marked relational structure is deployed in the dynamic, embodied,
context-sensitive activity of making sense of the world.

Active inference {[}@friston2017active{]} provides exactly this layer.
\end{frame}

\begin{frame}{Active Inference as a Process Theory of Language
Understanding}
\protect\phantomsection\label{active-inference-as-a-process-theory-of-language-understanding}
\begin{block}{The Free Energy Principle and Surprise Minimization}
\protect\phantomsection\label{the-free-energy-principle-and-surprise-minimization}
Active inference is the process theory derived from the free energy
principle (FEP): every self-organizing system maintains itself by
minimizing the surprisal (negative log-probability) of its sensory
observations under a \emph{generative model} of its environment
{[}-@friston2010free{]}. The system does this through two complementary
strategies:

\begin{enumerate}
\tightlist
\item
  \textbf{Perceptual inference}: Update internal beliefs to better
  predict current observations (reduce prediction error)
\item
  \textbf{Active inference}: Act on the environment to bring
  observations in line with predictions (reduce expected prediction
  error)
\end{enumerate}

Recent extensions of active inference to linguistics and cognitive
science have modeled language comprehension and production as forms of
sequential Bayesian inference. As Donnarumma, Frosolone, and Pezzulo
(2023) note in their integration of large language models and active
inference, ``linguistic processing {[}is{]} inference over a
hierarchical generative model, facilitating predictions and inferences
at various levels of granularity, from syllables to sentences''
{[}-@donnarumma2023integrating{]}. Similarly, Friston et al.~(2021) have
demonstrated how communication emerges between synthetic subjects:
``linguistic outcomes (specifically, the spoken word)\ldots{} are
selected to minimise the free energy given current beliefs'' via
``high-order interactions among abstract (discrete) states in deep
(hierarchical) models'' {[}-@friston2021understanding;
-@friston2020generative{]}.

Both strategies minimize the same quantity---variational free
energy---and both are driven by a single generative model that encodes
the system's expectations about the structure of its world.
\end{block}

\begin{block}{Language Understanding as Active Inference over Relational
Structure}
\protect\phantomsection\label{language-understanding-as-active-inference-over-relational-structure}
Language understanding on this view is \emph{active inference over
relational structure}: the listener maintains a generative model of
who-does-what-to-whom, and each incoming word provides evidence that
updates this model. Case marking provides crucial evidence---a
nominative suffix strongly predicts that the marked NP is the agent,
reducing uncertainty about the relational structure of the unfolding
event.

The process unfolds as follows:

\begin{enumerate}
\tightlist
\item
  \textbf{Prior}: The listener has a prior belief about the relational
  structure (a ``case diagram'' encoding expected roles and their
  connections)
\item
  \textbf{Observation}: Each word provides sensory evidence---its form,
  its case marking, its distributional properties
\item
  \textbf{Update}: The listener updates the case diagram to accommodate
  the evidence, using approximate Bayesian inference (typically
  variational message passing)
\item
  \textbf{Prediction}: The updated diagram generates predictions about
  upcoming words (case-marked NPs, verb valency patterns)
\item
  \textbf{Action}: In production, the speaker selects words and case
  markers that minimize expected free energy---choosing expressions that
  are informative, contextually appropriate, and syntactically
  well-formed
\end{enumerate}
\end{block}

\begin{block}{Connection to Barwise and Perry's Situation Semantics}
\protect\phantomsection\label{connection-to-barwise-and-perrys-situation-semantics}
This active inference perspective connects directly to the
\textbf{situation semantics} of Barwise and Perry
{[}-@barwise1983situations{]}, which conceptualized meaning as
structured \emph{situations}---collections of typed entities,
properties, and relations individuated by spatial and temporal location.
In our framework:

\begin{itemize}
\tightlist
\item
  A \textbf{situation} corresponds to an instantiated case diagram: a
  specific assignment of entities to case roles, with particular
  morphisms activated
\item
  The \textbf{situation type} corresponds to the case category itself:
  the abstract pattern of roles and relations that the situation
  instantiates
\item
  \textbf{Information flow} between situations corresponds to functorial
  mappings between case categories
\end{itemize}

Active inference adds the dynamic component: the agent moves through a
sequence of situations, updating its case diagram in real time and using
the diagram to predict which situation will come next.
\end{block}
\end{frame}

\begin{frame}{Distributional Active Inference: Convergence of Two
Distributional Traditions}
\protect\phantomsection\label{distributional-active-inference-convergence-of-two-distributional-traditions}
A remarkable convergence has recently emerged between
\emph{distributional semantics} in linguistics and \emph{distributional
reinforcement learning} in machine learning, mediated by active
inference. Akgül et al. {[}-@akgul2026distributional{]} introduce
\textbf{Distributional Active Inference (DAIF)}, which integrates active
inference into the distributional RL framework of Bellemare, Dabney, and
Munos {[}-@bellemare2017distributional{]}. Where classical RL optimizes
\emph{expected} returns (scalar values), distributional RL models the
\emph{full distribution} of returns---a shift from point estimates to
distributional representations that parallels the shift from symbolic to
distributional semantics in linguistics.

The formal architecture of DAIF proceeds through three stages: (1)
reconstructing active inference via variational Bayesian inference on a
controlled Markov process, expressing priors through Pearl's
do-calculus; (2) defining a \emph{push-forward} operation on
representation paths that maps latent-space trajectories to return
distributions; and (3) deriving a temporal-difference quantile-matching
algorithm that implements active inference without requiring explicit
transition dynamics modeling. The resulting ``push-forward RL'' template
achieves active inference's sample-efficiency advantages within a
model-free computational architecture:

\[\mathbb{E}\left[\sum_{t=0}^{\infty} \gamma^t R(x_t, a_t) \mid x_0, a_0\right] = \int_{\mathcal{S}^{\mathbb{N}_+}} R \circ f \, d(\mathbf{S}_{\#} \mathbb{P}_{x_0, a_0}^{P_\pi}) \tag{7.1}\]

where \(\mathbf{S}_{\#}\) denotes the push-forward measure on
representation paths and \(f: \mathcal{S} \to \mathcal{X}\) is the
stochastic decoder.

The terminological collision between ``distributional'' in
distributional semantics and ``distributional'' in distributional RL is
not mere homonymy---it reflects a deep structural parallel. In both
domains, the core move is the same: \textbf{replacing scalar summaries
with full distributional representations.} In linguistics, this means
replacing symbolic word identities with probability distributions over
contexts (Firth's {[}-@firth1957papers{]} company-keeping principle). In
RL, it means replacing expected-value estimates with full return
distributions. In active inference, it means replacing point estimates
of world states with variational posterior distributions. The
enriched-categorical framework of \autoref{sec:enriched-categories}
provides the unifying abstraction: all three are instances of
\([0,1]\)-enriched categories where hom-values encode distributional
proximity rather than discrete identity.

For case-theoretic reasoning, DAIF suggests a computational architecture
in which case assignment operates distributional-ly at every level: the
agent maintains not a single case diagram but a \emph{distribution over
case diagrams}, weighted by their posterior probability given the
observed linguistic evidence. Each incoming word updates this
distribution via variational message passing, and the agent's production
choices minimize expected free energy across the full distribution of
possible relational structures---not merely the most likely one. This
distributional perspective on case assignment aligns naturally with the
graded proto-role structure of Dowty {[}-@dowty1991thematic{]}: a noun
phrase is not categorically ``agent'' or ``patient'' but distributes
probability mass across case roles, with the distribution sharpening as
more evidence accumulates.
\end{frame}

\begin{frame}{CEREBRUM: The Computational Architecture}
\protect\phantomsection\label{cerebrum-the-computational-architecture}
\begin{block}{Architecture and Design Principles}
\protect\phantomsection\label{architecture-and-design-principles}
The \textbf{CEREBRUM} framework {[}@friedman2024cerebrum;
@cerebrum2024github{]}---Case-Enabled Reasoning Engine with Bayesian
Representations for Unified Modeling---provides a computational
architecture that implements the categorical case framework within an
active inference engine. CEREBRUM instantiates the view of Vasil et al.
{[}-@vasil2020world{]} that human communication is itself active
inference: a process of jointly constructing and refining generative
models of shared relational structure.

CEREBRUM's key design principles:

{\def\LTcaptype{none} % do not increment counter
\begin{longtable}[]{@{}
  >{\raggedright\arraybackslash}p{(\linewidth - 2\tabcolsep) * \real{0.5000}}
  >{\raggedright\arraybackslash}p{(\linewidth - 2\tabcolsep) * \real{0.5000}}@{}}
\toprule\noalign{}
\begin{minipage}[b]{\linewidth}\raggedright
Principle
\end{minipage} & \begin{minipage}[b]{\linewidth}\raggedright
Implementation
\end{minipage} \\
\midrule\noalign{}
\endhead
\textbf{Cases as functional roles} & Model components carry case
markings that determine their computational role in the inference
cycle \\
\textbf{Morphisms as message passing} & Grammatical relations are
implemented as message-passing channels between components \\
\textbf{Enriched weights as precision} & The \([0,1]\) weights on
morphisms correspond to precision parameters in the variational
inference scheme \\
\textbf{Alignment as model selection} & Different alignment types
correspond to different generative model architectures, selected by
Bayesian model comparison \\
\textbf{Diagrams as generative models} & Commutative diagrams serve as
the structural specification of the generative model \\
\bottomrule\noalign{}
\end{longtable}
}
\end{block}

\begin{block}{Case Roles as Functional Specializations in CEREBRUM}
\protect\phantomsection\label{case-roles-as-functional-specializations-in-cerebrum}
CEREBRUM deploys the eight traditional cases as functional
specializations:

{\def\LTcaptype{none} % do not increment counter
\begin{longtable}[]{@{}lll@{}}
\toprule\noalign{}
Case & CEREBRUM Function & Active Inference Role \\
\midrule\noalign{}
\endhead
NOM & Primary driver / agent & Source of action policies \\
ACC & Primary target / patient & Object of predictions \\
GEN & Source / possessor & Provider of priors \\
DAT & Recipient / goal & Target of information transfer \\
INS & Instrument / means & Tool for state transformation \\
LOC & Context / environment & Markov blanket boundary \\
ABL & Origin / cause & Source of causal influence \\
VOC & Addressee & Pragmatic pointer \\
\bottomrule\noalign{}
\end{longtable}
}
\end{block}
\end{frame}

\begin{frame}{Diagrammatic Reasoning as a Cognitive Process Theory}
\protect\phantomsection\label{diagrammatic-reasoning-as-a-cognitive-process-theory}
\begin{block}{Diagrams as the Format of Generative Models}
\protect\phantomsection\label{diagrams-as-the-format-of-generative-models}
The claim from \autoref{sec:introduction} can now be strengthened:
commutative diagrams are not merely convenient representations of case
structure---they are the \emph{format} in which cognitive agents
maintain their generative models of relational structure.

This claim is supported by converging evidence:

\begin{enumerate}
\item
  \textbf{Computational advantage} (Larkin \& Simon,
  {[}-@larkin1987diagram{]}): Diagrams enable search, recognition, and
  inference operations that are computationally prohibitive in
  sentential format. A commutative case diagram allows the agent to
  verify consistency (does the direct path equal the composed path?) by
  simple spatial inspection.
\item
  \textbf{Free ride inferences} (Shimojima,
  {[}-@shimojima1996reasoning{]}): Properties of the diagram that are
  perceptually available but would require explicit computation in a
  sentential format. In a case diagram, transitivity of grammatical
  relations is \emph{visible}---the existence of a path from NOM to DAT
  through ACC is spatially apparent.
\item
  \textbf{Hybrid reasoning} (Giardino,
  {[}-@giardino2017diagrammatic{]}): Mathematical diagrams engage a mode
  of reasoning that combines perceptual pattern recognition with
  background theoretical knowledge. Case diagrams similarly engage both
  the perceptual system (spatial layout) and linguistic knowledge (case
  constraints, verb valency).
\item
  \textbf{Peirce's existential graphs}: Peirce's graphical logic system
  demonstrated that first-order logic can be conducted entirely
  diagrammatically, without algebraic symbols. Our case diagrams extend
  this tradition: the relational structure of a sentence is represented
  graphically, and inference proceeds by diagram manipulation
  (adding/removing nodes, composing morphisms).
\end{enumerate}
\end{block}

\begin{block}{Predictive Processing and Diagrammatic Belief Updating}
\protect\phantomsection\label{predictive-processing-and-diagrammatic-belief-updating}
The predictive processing framework---of which active inference is the
most developed version---provides a natural account of how diagrammatic
representations are used cognitively:

\begin{enumerate}
\item
  \textbf{Top-down predictions}: The current case diagram generates
  predictions about expected sensory input (e.g., ``a nominative-marked
  NP should appear because the transitive verb requires an agent'')
\item
  \textbf{Bottom-up prediction errors}: Incoming words that violate the
  diagrammatic predictions generate prediction errors (e.g., an
  unexpected case marker triggers a P600 event-related potential in the
  brain)
\item
  \textbf{Belief updating}: The diagram is updated to accommodate the
  prediction error, potentially restructuring the assignment of entities
  to case roles (garden-path reanalysis)
\item
  \textbf{Precision weighting}: The enriched weights on morphisms serve
  as \emph{precision parameters} that control the relative influence of
  prior expectations and incoming evidence. A high-weight morphism
  generates strong predictions that are costly to override; a low-weight
  morphism generates weak predictions that are easily overridden.
\end{enumerate}
\end{block}
\end{frame}

\begin{frame}{Total Cognitive Scenario Understanding: The Integrated
Framework}
\protect\phantomsection\label{total-cognitive-scenario-understanding-the-integrated-framework}
The full picture emerges when we combine all five pillars within the
active inference framework:

\begin{enumerate}
\tightlist
\item
  \textbf{Case categories} (\autoref{sec:case-systems}) provide the
  \emph{objects and morphisms} of the generative model---the vocabulary
  of roles and relations
\item
  \textbf{Categorial grammar} (\autoref{sec:categorial-grammar})
  provides the \emph{composition rules}---how roles combine to form
  structured derivations
\item
  \textbf{DisCoCat/DisCoCirc} (\autoref{sec:categorical-semantics})
  provides the \emph{semantic functor}---mapping syntactic structure to
  distributional meaning
\item
  \textbf{Enriched structure} (\autoref{sec:enriched-categories})
  provides the \emph{precision parameters}---graded weights that control
  inference
\item
  \textbf{Topos-theoretic bridges} (\autoref{sec:topos-theory}) provide
  \emph{transfer theorems}---ensuring consistency across formalizations
\end{enumerate}

The active inference agent uses this combined structure as a single,
integrated generative model. Each scenario it encounters---a sentence
heard, a scene observed, an action planned---is interpreted by
instantiating a case diagram from structure (1), parsing the input using
rules (2), computing meaning via the semantic functor (3), weighting
confidence using enriched structure (4), and transferring results across
representational formats using bridge techniques (5).

This is \emph{total cognitive scenario understanding}: the agent doesn't
just parse a sentence or assign case labels---it constructs a complete,
internally consistent, generic, strongly typed, dynamically updating
model of the relational structure of the situation, and uses that model
to predict, explain, and act.
\end{frame}

\begin{frame}[fragile]{Computational Verification and Test Suite}
\protect\phantomsection\label{computational-verification-and-test-suite}
The framework developed in this paper is not merely theoretical---it is
computationally verified through a comprehensive implementation and test
suite that exercises every categorical construction discussed above.

\begin{block}{Implementation Architecture and Categorical Core}
\protect\phantomsection\label{implementation-architecture-and-categorical-core}
The categorical core (\texttt{CaseCategory}, \texttt{EnrichedCategory},
\texttt{CaseFunctor}, \texttt{NaturalTransformation}) is implemented in
Python with networkx-backed directed graphs, enforcing categorical
axioms at construction time. The visualization layer produces all 15
manuscript figures programmatically, ensuring exact correspondence
between formal claims and visual evidence. The DisCoPy integration
library (version 1.2.2) provides an independent validation path:
pregroup types (\texttt{Ty}), lexical entries (\texttt{Word}), cup
contractions (\texttt{Cup}), cap expansions (\texttt{Cap}), type
permutations (\texttt{Swap}), normal form computation
(\texttt{normal\_form()}), and circuit depth analysis (\texttt{depth()})
are exercised against the same categorical structures described in
\autoref{sec:categorial-grammar} and
\autoref{sec:categorical-semantics}.
\end{block}

\begin{block}{Test Suite and Coverage Metrics}
\protect\phantomsection\label{test-suite-and-coverage-metrics}
The implementation is validated by \textbf{158 automated tests} across 8
test files with \textbf{95.97\% code coverage} (against a 90\% threshold
enforced in the build configuration). Every test uses real mathematical
computations---no mocks or fakes:

\begin{itemize}
\tightlist
\item
  \textbf{Categorical axiom tests}: Identity morphism existence,
  composition associativity, weight invariants
\item
  \textbf{Enriched category tests}: Hom-value constraints, composition
  inequality, categorical magnitude
\item
  \textbf{Diagram type tests}: Pregroup diagrams validated for
  \texttt{dom\ ==\ Ty()} and \texttt{cod\ ==\ s}, correct box counts,
  diagram equality
\item
  \textbf{Metrics tests}: Normal form preservation, depth computation
  with graceful fallback for pregroup diagrams
\end{itemize}

This computational verification demonstrates that the category-theoretic
framework is not just a mathematical convenience but a \emph{working
computational architecture}---the categorical abstractions compile,
execute, and produce verifiable results, bridging the gap between formal
theory and implemented system.
\end{block}
\end{frame}

\end{document}
